\documentclass[12pt,a4paper]{article}
\usepackage[top=22mm,bottom=24mm,left=27mm,right=27mm]{geometry}
\usepackage{amsmath,amssymb}
\usepackage{xcolor}
\definecolor{link}{HTML}{205973}
\usepackage[colorlinks=true,linkcolor=link,urlcolor=link]{hyperref}
\hypersetup{pdftitle={Theorem 1. An approximation guarantee for the sampled value function}}
\setlength{\parindent}{0pt}
\setlength{\parskip}{5pt}
\raggedbottom
\newcommand{\R}{\mathbb R}
\newcommand{\lead}[1]{\par\smallskip\noindent\textbf{#1}\enspace}

\begin{document}
{\LARGE\bfseries Theorem 1. An approximation guarantee\\[2pt] for the sampled value function\par}
\medskip
Theorem~1 compares a value function fitted from sampled states with the
best approximation available in the same class. Its proof uses the
discounted state distribution both to choose the sample and to measure
one-step errors. The occupancy lemma below is verified in Lean. The full
probability bound remains open in the formalization.

\section*{The sampled problem and the target bound}
Let \(\mathcal S,\mathcal A\) be finite, nonempty state and action sets,
and let \(0<\alpha<1\) be the discount. Action \(a\) at state \(x\) costs
\(g_{x,a}\in\R\) and moves to \(y\) with probability \(p(x,y,a)\).
The Bellman operator for cost minimization is
\[
 (TJ)(x)=\min_{a\in\mathcal A}
   \left\{g_{x,a}+\alpha\sum_{y\in\mathcal S}p(x,y,a)J(y)\right\}.
\]
Write \(J^*\) for the optimal discounted cost-to-go and fix an optimal policy
\(\mu^*\). Let \(\nu\) be an initial probability distribution. For any
value vector \(f\), write \(\nu^\top f=\sum_x\nu_x f(x)\),
\(\|f\|_{1,\nu}=\sum_x\nu_x|f(x)|\), and \(\|f\|_\infty=\max_x|f(x)|\).

A feature map \(\Phi\) takes states into a real Hilbert space \(\mathcal H\).
For radii \(C,B\ge0\), the candidate functions are
\(\widetilde J_{\mathbf z,b}(x)=\langle\Phi(x),\mathbf z\rangle+b\),
with \(\|\mathbf z\|_{\mathcal H}\le C\) and \(|b|\le B\).
Draw \(X_1,\ldots,X_N\) independently from \(\pi_{\mu^*,\nu}\), the
discounted state distribution defined in \eqref{eq:occ}. Fit the parameters
by solving
\[
\begin{aligned}
 \operatorname{maximize}_{\mathbf z,b,s}\quad&
 \nu^\top\widetilde J_{\mathbf z,b}
       -\frac{2}{(1-\alpha)N}\sum_{i=1}^N s_i\\
 \text{subject to}\quad&
 \widetilde J_{\mathbf z,b}(X_i)\le (T\widetilde J_{\mathbf z,b})(X_i)+s_i,
 \quad s_i\ge0 \quad (1\le i\le N),\\
 &\|\mathbf z\|_{\mathcal H}\le C,\quad |b|\le B .
\end{aligned}
\tag{18}
\]
\textbf{Theorem 1 (target).} Let \(\epsilon>0\), \(0<\delta<1\), and
\(N\ge1\). With \(K=\max_x\|\Phi(x)\|_{\mathcal H}\) and
\(\|g\|_\infty=\max_{x,a}|g_{x,a}|\), set
\[
 \Xi=\bigl(4CK(1+\alpha)+4B(1-\alpha)+2\|g\|_\infty\bigr)
            \bigl(1+\sqrt{2\ln(1/\delta)}\bigr).
\]
If \(N\ge\Xi^2/\epsilon^2\), Theorem~1 asserts that, with probability
at least \(1-\delta-\delta^4\), an optimizer
\((\widehat{\mathbf z},\widehat b,\widehat s)\) of (18) satisfies
\[
 \bigl\|J^*-\widetilde J_{\widehat{\mathbf z},\widehat b}\bigr\|_{1,\nu}
 \le \frac{3+\alpha}{1-\alpha}
       \inf_{\|\mathbf z\|_{\mathcal H}\le C,\,|b|\le B}
       \bigl\|J^*-\widetilde J_{\mathbf z,b}\bigr\|_\infty
       +\frac{4\epsilon}{1-\alpha}.
 \tag{19}
\]

\newpage
\section*{Discounted occupancy and the Bellman residual}
Fix any policy \(\mu\), and set \(g_\mu(x)=g_{x,\mu(x)}\) and
\(P_\mu(x,y)=p(x,y,\mu(x))\). Its expected cost vector and its one-step
update of a candidate \(J\) are
\[
 J_\mu=\sum_{t=0}^\infty\alpha^tP_\mu^t g_\mu,\qquad
 T_\mu J=g_\mu+\alpha P_\mu J ,
\]
where \(P_\mu^0=I\), the identity matrix. The residual \(J-T_\mu J\)
measures how far the proposed cost exceeds its one-step prediction.

\textbf{Lemma (discounted occupancy).} For every \(\mu\) and \(\nu\) as above, the row vector
\begin{align}
 \pi_{\mu,\nu}^{\top}
 =(1-\alpha)\sum_{t=0}^{\infty}\alpha^t\nu^\top P_\mu^t
 =(1-\alpha)\nu^\top(I-\alpha P_\mu)^{-1}
 \tag{17}\label{eq:occ}
\end{align}
is a probability distribution. It is the unique real row vector satisfying
\begin{align}
 \pi_{\mu,\nu}^{\top}=(1-\alpha)\nu^\top+\alpha\pi_{\mu,\nu}^{\top}P_\mu .
 \tag{B}\label{eq:balance}
\end{align}
For every \(J\in\R^{\mathcal S}\), the expected initial value error satisfies
\begin{align}
 \nu^\top(J-J_\mu)
   =\frac{1}{1-\alpha}\,\pi_{\mu,\nu}^{\top}(J-T_\mu J).
 \tag{R}\label{eq:residual}
\end{align}

\lead{Proof. The distribution of discounted visits.}
The row vector \(\nu^\top P_\mu^t\) is the state distribution at time \(t\).
Indeed, multiplying a probability row by \(P_\mu\) preserves
nonnegativity and total mass. The coefficients \((1-\alpha)\alpha^t\)
are nonnegative and sum to one. Thus their mixture in \eqref{eq:occ}
converges coordinatewise and has total mass one. A state with zero initial
mass can have positive occupancy if the policy can reach it later.

\lead{The inverse and the balance equation.}
Every entry of \(P_\mu^t\) lies between zero and one, since each row is
a probability distribution. The matrix series \(\sum_{t\ge0}\alpha^tP_\mu^t\)
therefore converges absolutely in each entry. For the partial sum
\(D_n=\sum_{t=0}^{n}\alpha^tP_\mu^t\), multiplication on either side gives
\[
  (I-\alpha P_\mu)D_n
      =D_n(I-\alpha P_\mu)=I-\alpha^{n+1}P_\mu^{n+1}.
\]
The last term on the right tends to zero entrywise. There are only finitely
many states, so multiplication by \(I-\alpha P_\mu\) commutes with the
limit. The sum is consequently its two-sided inverse. This proves the
second formula in \eqref{eq:occ}.

\newpage
\section*{Balance, uniqueness and the residual}
Multiplying \eqref{eq:occ} on the right by \(I-\alpha P_\mu\) gives
\[
   \pi_{\mu,\nu}^\top(I-\alpha P_\mu)=(1-\alpha)\nu^\top.
\]
Rearranging gives \eqref{eq:balance}. Conversely, any real row vector
satisfying that balance equation must give the same product with
\(I-\alpha P_\mu\). Multiplication by the inverse determines the row
uniquely. In \eqref{eq:balance}, the first term accounts for time zero
and the second for visits after one or more transitions.

\lead{From a one-step error to the initial error.}
The series for \(J_\mu\) gives \(J_\mu=(I-\alpha P_\mu)^{-1}g_\mu\).
Hence \(J_\mu\) is the unique solution of \(J_\mu=g_\mu+\alpha P_\mu J_\mu\).
Subtracting this equality from the definition of the residual gives
\[
 (I-\alpha P_\mu)(J-J_\mu)
     =J-\alpha P_\mu J-g_\mu=J-T_\mu J .
\]
Multiply on the left by \(\pi_{\mu,\nu}^\top\) and use the balance equation.
The left side becomes \((1-\alpha)\nu^\top(J-J_\mu)\). Division by
\(1-\alpha>0\) gives \eqref{eq:residual}. \hfill\(\square\)

\section*{How the lemma enters Theorem 1}
For the optimal policy, \(J_{\mu^*}=J^*\). Write
\(\Delta^*=(I-\alpha P_{\mu^*})^{-1}\) and let
\(s_{\mu^*}(x)=\max\{J(x)-(T_{\mu^*}J)(x),0\}\) be the positive
one-step residual. Applied to this vector, \eqref{eq:occ} gives
\[
 \nu^\top\Delta^*s_{\mu^*}
       =\frac{1}{1-\alpha}\pi_{\mu^*,\nu}^\top s_{\mu^*}.
\]
The left side weights the accumulated residual by the starting distribution.
The right side averages a one-step residual under the law used to draw the
states in (18). This is the substitution in the paper's inequality chain (26).
The next deterministic step bounds overestimation by
\(\Delta^*s_{\mu^*}\). The sampling estimates must then control the
difference between the population and empirical residual means.

\newpage
\section*{A two-state example}
The states are \(s,t\), and the initial state is \(s\). From either state,
action \(A\) costs \(1\) and moves to \(s\) or \(t\) with equal probability.
Action \(B\) costs \(2\) and moves to \(t\) with certainty. The policy
takes \(A\) at \(s\) and \(B\) at \(t\). In state order \(s,t\),
\[
 P_\mu=\begin{pmatrix}1/2&1/2\\0&1\end{pmatrix},\qquad
 g_\mu=\begin{pmatrix}1\\2\end{pmatrix},\qquad
 \nu^\top=(1,0),\qquad \alpha=\tfrac12 .
\]
Writing \(u=\pi_{\mu,\nu}(s)\), the first coordinate of
\eqref{eq:balance} is \(u=1/2+u/4\). Thus the occupancy is
\((2/3,1/3)\). State \(t\) accounts for a third of discounted time even
though the initial distribution gives it zero mass.

At \(t\), the policy pays \(2\) forever, giving \(J_\mu(t)=4\). At \(s\),
the Bellman equation gives
\(J_\mu(s)=1+\tfrac14 J_\mu(s)+\tfrac14 J_\mu(t)\), so \(J_\mu(s)=8/3\).
Take the candidate \(J(s)=3,\ J(t)=2\). Its one-step predictions and
errors are
\[
 \begin{array}{c@{\qquad}r@{\quad}r@{\quad}r@{\quad}r}
   \text{State} & J_\mu & J & T_\mu J & J-T_\mu J\\ \hline
   s & 8/3 & 3 & 9/4 & \phantom{-}3/4\\
   t & 4   & 2 & 3   & -1
 \end{array}
\]
The residuals at the two states have opposite signs. Their contributions
partly cancel when averaged over discounted visits,
\[
 \pi_{\mu,\nu}^\top(J-T_\mu J)=\frac23\cdot\frac34+\frac13\cdot(-1)=\frac16 .
\]
Dividing by \(1-\alpha=1/2\) gives an initial error of \(1/3\). Direct
evaluation at the starting state gives the same answer,
\(J(s)-J_\mu(s)=3-8/3=1/3\).

\section*{Relation to the original result}
The paper introduces (17) with the optimal policy in (18). Its formula holds
as printed, and the lemma proves it for every fixed policy. Appendix~B
obtains Theorem~1 from the weighted Theorem~2 by setting its weight
function \(\psi\) equal to one. The target on page~1 uses the confidence
range \(0<\delta<1\) so the square root in \(\Xi\) is real. The paper's
statement allows \(\delta>0\).

\smallskip
{\small
\textbf{Source.} N.~Bhat, V.~F.~Farias, C.~C.~Moallemi and A.~T.~Zheng,
\href{https://moallemi.com/ciamac/papers/kernel-adp-2012.pdf\#page=12}{\emph{Non-parametric Approximate Dynamic Programming via the Kernel Method}},
August 29, 2022, \S\S3.2--3.3 and Appendix~B.\par
\textbf{Formal proof.}
\href{https://github.com/nikhilpb/kadp/blob/main/Kadp/Occupancy.lean}{Lean occupancy and residual identities}
and \href{https://github.com/nikhilpb/kadp/blob/main/Kadp/Examples/Occupancy.lean}{two-state example}.
The series proof applies
\href{https://github.com/LeanEconomics/DynamicProgramming/blob/9d4783f3fb282444f8f3a91c4e6439293dd4bacc/FiniteStates/MarkovDecisionProcesses/MarkovDecisionProcesses/MDP.lean}{Robert Kirkby's policy-evaluation theorems}
to each state's indicator function.}
\end{document}
